% Immutable cited-input supplement. Target proofs remain embargoed.
% Target: cdp:st:C.2.D4.2; tag 01BX.
% Frozen source SHA256 74497bfe327fae25c9cdcb17334b3bf42a9793a9df7c0a4404d97568f2f760f0.
% CompendiumRef is resolved against the cited S6_09 source edition, not silently identified with the current series statement.

% Exact cited input: series:II / cdp:st:A.5.19
\begin{proposition}[Ordinary-triple images have arithmetic genus at most one]\hypertarget{tag-01B9}{}\label{cdp:st:A.5.19}


Let $f:Y\to C_0$ be a proper holomorphic map from a smooth
compact complex threefold with $b_2(Y)=0$, with connected
fibres and smooth general fibre a complex two-torus.
Let $mS$ be a full fibre with irreducible reduced support,
and let $D\subset\operatorname{Sing}S$ be an irreducible
curve with generic ordinary-triple plane transversal.
Then $D$ is a local complete intersection in $Y$ and
$p_a(D)\le1$. The possibilities for its underlying curve
are a smooth rational curve, a rational curve with one
ordinary node or cusp, and a smooth elliptic curve.
In the elliptic case the transversal type is constant.

When $D$ is singular rational, every component of the
normalized conductor curve dominating $D$ is rational,
and its map to $D^\nu\simeq\PP^1$ has degree one.
Consequently, in the irrational ruled setup of \ref{cdp:st:A.5.18}
with the generic transversal replaced by an ordinary
triple point, the image of the noncontracted horizontal
section is either smooth rational or smooth elliptic.
In the elliptic case the horizontal cover is unramified
of degree three, the normalization is smooth minimal
elliptic ruled, and the fibre has Euler number zero.
In the rational case the horizontal cover has degree
two or three; these cases are not excluded by this
statement.
\end{proposition}

% Exact cited input: comp:II / cdp:st:C.2.RAT.6
\begin{proposition}[Credited re-derivation: Cubic suspensions of plane curves with delta at most five]\label{cdp:st:C.2.RAT.6}
Let $r(x,y)=0$ be a reduced plane-curve germ with
$\delta(r)\le5$, and let $X_r$ be the isolated surface
singularity $z^3=r(x,y)$. Then $p_g(X_r)\le1$ and the
first rational Betti number of its link is at most two.
The latter is two precisely for plane type $A_5$ or $D_4$.
For a collection of these plane germs with total delta
invariant at most five,
\[
             \sum p_g(X_r)\le1,
           \qquad \sum b_1(\operatorname{link}X_r)\le2.
\tag{C.RD.5}\label{cdp:eq:C.RD.5}
\]
\end{proposition}

% Exact cited input: series:II / cdp:st:A.5.18
\begin{proposition}[Every horizontal tacnode image is elliptic]\hypertarget{tag-01B8}{}\label{cdp:st:A.5.18}


Let $f:Y\to C_0$ be a proper holomorphic map from a smooth
compact complex threefold with $b_2(Y)=0$, with connected
fibres and smooth general fibre a complex two-torus.
Let $mS$ be a full fibre with irreducible reduced support,
$\eta:T\to S$ its normalization, and $\pi:Z\to T$ a
resolution with an irrational ruling $p:Z\to B$.
Suppose every exceptional curve of $\pi$ is vertical and
the conductor divisor on $Z$ is $2C+V$, where $C$ is a
section of $p$ and $V$ is vertical. Let $D=\eta\pi(C)$.
Suppose the generic plane-curve transversal to $D$ in $S$ is an
ordinary tacnode.

Then $D$ is a smooth elliptic curve, $B=C$ is elliptic,
and $C\to D$ is an unramified double cover. The transversal
type along $D$ is everywhere the same ordinary tacnode.
The normalization $T$ is a smooth minimal ruled surface
over this elliptic curve. Its conductor is twice a
section, and $e(S)=0$. In particular its minimal
resolution is the identity.
In particular a rational image of
the horizontal tacnodal conductor cannot occur. The proof
allows vertical conductor and rational singularities of
$T$; it does not assume that the normalization map is
immersive at the special points under examination.
\end{proposition}

